\chapter{CJK Mathematics}
\label{ch:cjk-math}

OmniLaTeX supports CJK (Chinese, Japanese, Korean) characters in
mathematical contexts through \texttt{luatexja}.  This chapter covers
math-specific CJK features; for full language support including
document-level configuration, see Chapter~\ref{ch:cjk}.

% =========================================================================
\section{CJK Math Fonts}
% =========================================================================

When using LuaLaTeX, \texttt{luatexja} provides native support for CJK
characters inside math mode.  CJK glyphs are automatically rendered
alongside Latin math symbols without additional configuration.

\begin{verbatim}
\[
  E = mc^2 \quad \text{能量} = mc^2
\]
\end{verbatim}

\texttt{luatexja} maps CJK characters to appropriate math fonts based on
the document's font configuration.  The main text CJK font is used by
default for CJK characters in math mode, ensuring visual consistency
between inline text and displayed equations.

For upright CJK text within math mode, use the standard \cs{text\{\}}
command:

\begin{verbatim}
$T_{\text{効率}} = 0.85$
% → T_効率 = 0.85
\end{verbatim}

% =========================================================================
\section{Ruby Annotations}
% =========================================================================

Ruby annotations (furigana in Japanese, pinyin in Chinese) place small
reading aids above CJK characters.  OmniLaTeX provides these via the
\texttt{luatexja-ruby} package.

\begin{verbatim}
\ruby{漢字}{かんじ}        % Japanese furigana
\ruby{汉字}{hànzì}          % Chinese pinyin
\end{verbatim}

Shorthand aliases are available for convenience:

\begin{table}[htbp]
  \centering
  \caption{Ruby annotation commands.}
  \label{tab:ruby-commands}
  \begin{tabular}{@{}lll@{}}
    \toprule
    Command & Alias & Description \\
    \midrule
    \texttt{\textbackslash ruby\{漢字\}\{かんじ\}} & -- & Full ruby annotation \\
    \texttt{\textbackslash furigana\{漢字\}\{かんじ\}} & \texttt{\textbackslash ruby} & Japanese furigana alias \\
    \texttt{\textbackslash pinyin\{汉字\}\{hànzì\}} & \texttt{\textbackslash ruby} & Chinese pinyin alias \\
    \bottomrule
  \end{tabular}
\end{table}

The annotation size can be adjusted globally:

\begin{verbatim}
\setRubyScale{0.5}   % Scale ruby text to 50% of base size
\end{verbatim}

Ruby annotations work in both text and math mode, making them useful for
annotating technical terms that contain CJK characters in equations.

% =========================================================================
\section{Vertical Writing}
% =========================================================================

\texttt{luatexja} supports vertical (top-to-bottom, right-to-left) text
layout through the \texttt{vertical} environment:

\begin{verbatim}
\begin{vertical}
  この文章は縦書きです。
  \[ E = mc^2 \]
\end{vertical}
\end{verbatim}

For inline vertical text spans, use:

\begin{verbatim}
\verticaltext{縦書きテキスト}
\end{verbatim}

Math environments within vertical mode are rotated so that equations
remain readable.  Displayed equations appear rotated 90 degrees
counter-clockwise relative to the vertical text flow, following Japanese
typesetting conventions (JIS~X~4051).

Note that vertical writing mode interacts with page layout settings.
Margins, floats, and captions are automatically adjusted by
\texttt{luatexja} when the vertical environment is active.

% =========================================================================
\section{CJK-Latin Spacing}
% =========================================================================

Proper inter-character spacing between CJK and Latin scripts is handled
by \texttt{luatexja}'s spacing engine.  A small amount of space is
automatically inserted at CJK--Latin boundaries to prevent characters
from appearing cramped.

\begin{verbatim}
\cjklatinspacing  % Enable CJK-Latin inter-character spacing
\end{verbatim}

Fine-grained control is available through:

\begin{verbatim}
\setCJKLatinSpacing{0.25em}  % Set explicit CJK-Latin gap
\end{verbatim}

The default spacing is determined by the font metrics and generally does
not require manual adjustment.  However, in math-heavy documents with
mixed CJK and Latin symbols, explicit spacing control can improve
readability when the default gap is too small or too large for the
chosen font combination.

For complete CJK language setup including font selection, input
methods, and document-level options, refer to Chapter~\ref{ch:cjk}.
